\section{How it works}
\label{sec:short:how}

One pass through the three layers, with no formalism. The full paper specifies each of these to the
depth at which it could be reimplemented; this section aims only at the depth at which it can be
understood.

\subsection{The chain: a rotation, not a race}

Every collateralised node is in a queue. Producing a block moves a node from the front of the queue
to the back, and the block reward goes to whoever produced it, split by tier according to a fixed
schedule. There is no difficulty to grind and nothing to mine; a node's expected income is a function
of how many nodes are in its tier and how long the rotation takes to come round.

Because every block carries the same fixed weight, the chain does not have a longest-chain rule in
the Bitcoin sense. Fork choice is a block-count race resolved by which block was seen first, and
reorganisations are capped at forty blocks by consensus. That is a deliberate simplification, and it
has a cost the full paper treats at length: with no cumulative-work signal in the header, a
lightweight client that verifies without a full node cannot be built on today's format.

\subsection{Where new \flux{} comes from}

Every block mints new \flux{}, and today that issuance is the whole of an operator's income. The
rate is \num{14.0}~\flux{} per block, split across the three tiers and a development fund by fixed
shares. It falls by ten percent each October --- to \num{12.6} at block \num{3071200} in 2026, then
\num{11.34}, and so on for twenty annual steps --- and then \emph{freezes} rather than continuing
toward zero: from the twentieth reduction onward the chain issues a constant tail of roughly
\num{1789219}~\flux{} a year. Total issued supply reaches \num{548043625.86}~\flux{} by 2045 and
crosses \num{560000000} around 2052.

Two things follow for a holder. There is no hard cap that issuance approaches and stops at --- the
tail is perpetual by design. And the emission curve is what v9 is built to \emph{replace} as the
source of operator income: at PNR activation the operators' share of new issuance is planned to be
cut in half, with revenue from applications taking its place (\Cref{sec:short:changes}).

\subsection{The cloud: publish, claim, run}

A deployment starts as a specification --- the container image, the resources it needs, how many
copies should run, how long it is paid for. The owner signs it, pays for it, and broadcasts it.
That is the whole act of deploying: there is no console and no account.

Nodes then \emph{claim} it. Every eligible node sees the specification, decides whether it can
host it, and if so announces its intention and waits ninety seconds. If nobody else has claimed the
same instance in that window, it installs and starts the container. If the announced count is
already met, it stands down. The application ends up running on several nodes at once, chosen by
whoever was fast rather than by any scheduler.

This is unusual, and the paper says so plainly. There is no placement algorithm in the sense a
cloud engineer would expect --- no bin-packing, no affinity, no capacity model. What there is
instead is an incentive: nodes want to be full, applications want to be replicated, and the ninety-
second collision rule keeps them from tripping over each other. It works at the current scale. What
it does not do is give a tenant any say over \emph{where} the application lands.

\emph{How you reach it.} A running application gets a name under \texttt{app.runonflux.io} (and
\texttt{app2.runonflux.io}), resolved by the network's own DNS and fronted by load balancers that
health-check each instance and route around the ones that fail. That layer is what makes an
application on several anonymous machines look like one service --- and it is operated by the
Foundation on a sixteen-host fleet, which \Cref{sec:short:limits} returns to.

\begin{figure}[H]
\centering
\begin{tikzpicture}[font=\small, node distance=9mm and 9mm,
  st/.style={draw=FluxBlue, fill=FluxBlue!8, rounded corners=2pt, minimum height=11mm,
             text width=27mm, align=center, font=\scriptsize},
  ar/.style={->, thick, FluxSlate}]
\node[st] (spec) {owner signs and pays for a specification};
\node[st, right=of spec] (bcast) {specification reaches every node};
\node[st, right=of bcast] (claim) {eligible nodes announce a claim, wait \SI{90}{\second}};
\node[st, right=of claim] (run) {uncontested claimants install and run};
\draw[ar] (spec) -- (bcast); \draw[ar] (bcast) -- (claim); \draw[ar] (claim) -- (run);
\node[below=3mm of run, font=\scriptsize, text=FluxSlate, text width=30mm, align=center]
  {runs on several nodes; expires when the paid period ends};
\end{tikzpicture}
\caption{From a signed specification to a running application. No account is opened at any step,
and no scheduler decides placement: nodes claim work, and a fixed wait resolves collisions.}
\label{fig:short:lifecycle}
\end{figure}

\subsection{Paying: a transaction, not a relationship}

Today, paying for a deployment means one thing: a single signed transaction, for the full price of
the period, to an address the chain knows. The price is a function of the resources reserved and the
duration --- CPU, memory, disk, instances, blocks --- and it is a \emph{reservation} price. You pay
for what you asked to hold, not for what you used. There is no metering, and therefore nothing to
meter, dispute, or attest.

When the paid period ends, the application stops. Renewing it is the owner's act, done the same way.
The protocol offers no subscriptions and no automatic renewal; operators who want that convenience
can buy it from a third-party service that holds a balance on their behalf, and the paper is
explicit that such a service is a business built on \Flux{}, not a feature of it.

\begin{figure}[H]
\centering
\begin{tikzpicture}[font=\small,
  box/.style={draw=FluxSlate, fill=FluxSlate!8, rounded corners=2pt, minimum height=10mm,
              text width=30mm, align=center, font=\scriptsize},
  hi/.style={box, draw=FluxBlue, fill=FluxBlue!10},
  ar/.style={->, thick, FluxSlate}, lbl/.style={font=\scriptsize, text=FluxSlate}]
\node[lbl, anchor=west] at (0,2.55) {\textbf{Today}};
\node[box] (t1) at (1.7,1.7) {owner pays the full price};
\node[box] (t2) at (6.2,1.7) {payment address};
\node[box] (t3) at (10.7,1.7) {operators are paid by \emph{new issuance}, by tier rotation};
\draw[ar] (t1) -- (t2);
\draw[ar, dashed, FluxSlate!60] (t2) -- node[above, lbl] {no link} (t3);

\node[lbl, anchor=west] at (0,0.35) {\textbf{After PNR (July 2027)}};
\node[hi] (p1) at (1.7,-0.5) {owner pays the full price};
\node[hi] (p2) at (6.2,-0.5) {consensus splits it: \SI{80}{\percent} / \SI{20}{\percent}};
\node[hi] (p3) at (10.7,-0.5) {the \SI{20}{\percent} feeds a pool that pays operators};
\draw[ar] (p1) -- (p2); \draw[ar] (p2) -- (p3);
\end{tikzpicture}
\caption{Where a payment goes. Today the money a tenant pays and the money an operator earns are
unconnected --- operators are paid by issuance. Under PNR, part of every payment flows into a pool
that pays operators instead. This is the hinge between the network as it is and as it is being
rebuilt.}
\label{fig:short:payment}
\end{figure}

\limitnote{The payment address is a single consensus-known address, and the split is applied by
consensus, not by the node software --- so a software release cannot change the distribution. What
the protocol does \emph{not} give a tenant is any remedy if a node fails to serve what was paid for.
Reservation pricing has no delivered-uptime quantity to compute a refund against, and the paper
records the decision that any such guarantee lives in the service layer, not the protocol.}
